% =====================================================================
%  2bat-ud6-4.tex  —  SESSIÓ 4: De la mostra a la població: la lògica del mostreig
%  (sense preàmbul: s'inclou des de main.tex)
% =====================================================================
\horatitol{Sessió 4 — De la mostra a la població: la lògica del mostreig}%
{Formalitzar, sense tecnicisme, per què una mostra a l'atzar i prou gran representa el tot, i per què tota inferència té incertesa.}

\subsection{Idea clau}

Aquesta sessió explica \textbf{per què funciona} la inferència. La idea, intuïtiva i sense
fórmules, és que una mostra triada \textbf{a l'atzar} i \textbf{prou gran} tendeix a
assemblar-se a la població, de manera que el que hi observem és una bona
\textbf{estimació} del que passa al tot. Però sempre amb un grau d'\textbf{incertesa}: mai
és exacta.

\begin{idea}
\textbf{Les dues condicions perquè una mostra representi el tot}
\begin{itemize}[leftmargin=1.4em, itemsep=2pt]
  \item \textbf{A l'atzar:} cada membre de la població ha de tenir opció de sortir triat;
        així s'evita afavorir cap perfil (s'evita el biaix).
  \item \textbf{Prou gran:} com més gran és la mostra, més s'acosta el resultat al valor
        real de la població.
\end{itemize}
\textbf{Connexió amb la UD5.} Si repetíssim el mostreig moltes vegades, els resultats es
distribuirien al voltant del valor real de la població \textbf{seguint una campana}
(normal). Això és el que permet posar \textbf{garanties} i \textbf{marges d'error} a les
conclusions.
\smallskip

\noindent\textbf{El revers, amb honestedat:} tota inferència té incertesa, perquè una
mostra mai és exactament igual que la població. Per això les conclusions s'expressen
\textbf{sempre amb un marge}, mai com a certeses absolutes.
\end{idea}

\begin{center}
\begin{tikzpicture}
\begin{axis}[udaxis, width=11cm, height=4.6cm,
    xlabel={\small resultat de la mostra}, ylabel={}, ytick=\empty,
    xmin=-4, xmax=4, ymin=0, ymax=0.45, xtick={0}, xticklabels={valor real}]
  \addplot[udblue, thick, domain=-4:4, samples=120]{0.3989*exp(-x^2/2)};
  \addplot[udblue!18, domain=-2:2, samples=80] {0.3989*exp(-x^2/2)} \closedcycle;
  \node[udnavy, font=\scriptsize] at (axis cs:0,0.16){la majoria de mostres};
  \node[udgrey, font=\scriptsize] at (axis cs:-3,0.06){poques};
  \node[udgrey, font=\scriptsize] at (axis cs:3,0.06){poques};
\end{axis}
\end{tikzpicture}

\figcap{Si es repetís el mostreig moltes vegades, els resultats caurien en forma de campana al voltant del valor real: la majoria a prop, poques lluny.}
\end{center}

\subsection{Exemples}

\begin{exemple}[per què l'atzar és tan important]
Triar la mostra \textbf{a l'atzar} (per exemple, per sorteig d'un llistat complet) fa que
no s'afavoreixi cap grup: tothom té opció de sortir. Si, en canvi, es tria «els primers que
passin per aquí», s'afavoreix un perfil concret i apareix el biaix. L'atzar és el que
garanteix que la mostra sigui un \emph{reflex reduït} de la població.
\end{exemple}

\begin{exemple}[per què la mida importa]
Amb una mostra molt petita (per exemple, $10$ persones), el resultat pot variar molt per
pura sort: n'hi ha prou que un parell de respostes siguin extremes per moure'l. Amb una
mostra gran (per exemple, $1\,000$), les respostes extremes es compensen i el resultat
s'estabilitza a prop del valor real. Més mida $\to$ menys incertesa.
\end{exemple}

\subsection{Exercicis resolts}

\begin{exresolt}[valorar la fiabilitat d'un mostreig]
Un estudi tria la mostra \textbf{per sorteig} d'un cens complet i enquesta $1\,200$
persones. Un altre enquesta les \textbf{primeres 1\,200} persones que passen per una plaça.
Quin és més fiable i per què?
\solucio
El primer és més fiable: la tria \textbf{a l'atzar} d'un cens complet fa que la mostra
representi la població sense afavorir cap perfil. El segon, tot i tenir la mateixa mida,
té \textbf{biaix de selecció} (només qui passa per aquella plaça, en aquell moment), així
que la mida gran no el salva.
\resposta{El primer (mostra a l'atzar); la mida no corregeix el biaix del segon.}
\end{exresolt}

\begin{exresolt}[expressar la incertesa]
Un estudi ben fet conclou que «el $40\%$ de la població opina X». Per què seria incorrecte
dir-ho com una certesa absoluta?
\solucio
Perquè el $40\%$ prové d'una \textbf{mostra}, no de tota la població: una altra mostra
donaria un valor una mica diferent. El resultat real de la població estarà \emph{a prop}
del $40\%$, però no exactament. Per això s'ha d'expressar amb un \textbf{marge}
(per exemple, $40\%\pm 3\%$), no com un valor exacte i segur.
\resposta{Perquè és una estimació d'una mostra; s'ha d'expressar amb un marge, no com a
certesa.}
\end{exresolt}

\subsection{Exercicis per fer}

\begin{exercicis}
\begin{exlist}
  \item Explica les dues condicions perquè una mostra representi bé la població.
  \item Per què una mostra triada «a mà» (per exemple, coneguts de l'enquestador) no és
        fiable, encara que sigui gran?
  \item Dos estudis mesuren el mateix, un amb $50$ persones i l'altre amb $1\,500$, tots
        dos a l'atzar. Quin donarà una estimació més precisa (menys incertesa)? Per què?
  \item Per què cap estudi seriós expressa el seu resultat com un valor exacte i segur?
  \item \textbf{(Connexió UD5)} Explica, amb la idea de la campana, per què repetir el
        mostreig moltes vegades dona resultats concentrats a prop del valor real.
  \item \textbf{(Raonament)} «Com més gran és la mostra, més s'acosta al valor real.» Vol
        dir això que amb una mostra molt gran ja no hi ha gens d'incertesa? Raona-ho.
\end{exlist}
\end{exercicis}

% ---------------------------------------------------------------------
\fullsolucions{Sessió 4}
\begin{solucions}
\begin{sollist}
  \item (1) A l'atzar: tothom de la població ha de tenir opció de sortir triat, per no
        afavorir cap perfil. (2) Prou gran: com més gran és la mostra, més s'acosta el
        resultat al valor real.
  \item Perquè trien un perfil concret (biaix de selecció): els coneguts s'assemblen entre
        ells i no reflecteixen la diversitat de la població; la mida no corregeix aquest
        biaix.
  \item El de $1\,500$: amb més mida les respostes extremes es compensen i el resultat
        s'estabilitza a prop del valor real (menys incertesa).
  \item Perquè el resultat prové d'una mostra i sempre té un grau d'incertesa; expressar-lo
        com a exacte amagaria aquesta incertesa i podria induir a error.
  \item Resposta oberta. Idea: si es repetís el mostreig moltes vegades, els resultats es
        distribuirien en forma de campana centrada al valor real; la majoria de mostres
        cauen a prop del centre i poques lluny, de manera que una mostra concreta sol
        estar a prop del valor real.
  \item Resposta oberta. Idea: no; la incertesa \textbf{disminueix} amb la mida, però mai
        arriba a zero mentre no es mesuri tota la població. Sempre queda un marge, encara
        que sigui petit.
\end{sollist}
\end{solucions}
